\documentclass[10pt]{article}
\usepackage[utf8]{inputenc}
\usepackage[T1]{fontenc}
\usepackage{amsmath}
\usepackage{amsfonts}
\usepackage{amssymb}
\usepackage[version=4]{mhchem}
\usepackage{stmaryrd}

\begin{document}
создающими эффект «овыпукления» (см. [6], [7]). Отсутствие выпуклости у интегральных функционалов $J(x(\cdot), u(\cdot))$ возмещают путем погружения задачи в более общую, с новым функционалом, являющимся выпуклой минорантой прежнего, и решения новой задачи в более широком классе управлений (см. [8]). В отмеченных случаях существование оптимального управления часто вытекает из существования допустимого управления.

Теория необходимых условий экстремума наиболее развита в задачах О. у. П. Основополагающим результатом здесь послужил принцип максимума Понтрягина, содержащий необходимые условия сильного экстремума в задаче оптимального управления.

Следует отметить создание общих приемов получения необходимых условий в экстремальных задачах, әффективно использованных для задач О.у. п. с более сложными ограничениями (фазовыми, функциональными, минимаксными, смешанными и т. д.) и основанных, так или иначе, на теоремах об отделимости выпуклых конусов (см. [9], [10]). Пусть, напр., $E-$ векторное пространство, $f(x), \quad x \in E,-$ заданный функционал, $Q_{i}$ - множества из $E$,

\[
Q=\cap Q_{i}, i=1, \ldots, n,
\]

$x^{0} \in Q$ - точка, в к-рой $f(x)$ достигает минимума на $Q$, $Q_{0}=\left\{x \mid f(x)<f\left(x^{0}\right)\right\}$

Суть распространенного общего метода состоит в том, что каждое из множеств $Q_{i}, i=0,1, \ldots, n$, аппроксимируется в окрестности точки $x^{0}$ нек-рым выпуклым конусом $K_{i}$ с вершиной в точке $x^{0}$ (конус "убывания" для $Q_{0}$; конус «допустимых направлений» для ограничениіи, изображаемых неравенствами; конус «касательных нашравлений» для ограничений типа равенства, в том числе для дифференциальных связей, и т. д.). Необходимое условие минимума теперь состоит в том, чтобы $x^{0}$ была единственной точкой, общей для всех $K_{i}, i=0,1, \ldots, n$, и, следовательно, чтобы конусы были «отделимы» (см. [8]). Последнему «геометрическому» условию далее придают аналитич. форму и по возможности преобразуют к удобному виду, напр. при помощи функции Гамильтона. В зависимости от исходных ограничениӥ, а также от класса используемых вариаций, необходимые условия могут принимать как форму, аналогичную принципу Понтрягина, так и форму локального (линеаризованного) принцигга максимума (условия слабого экстремума по $u$ ). Реализация отмеченного пути, таким образом, зависит от возможности аналитически описывать конусы $K_{i}$. Әффективное их описание достигается для множеств $Q_{i}$, задаваемых гладкими функциями, удовлетворяющими нек-рым дополнительным условиям регулярности в рассматриваемой точке, или выпуклыми функциями (см. [9], [10]). В принципе отмеченный путь допускает обобщения и на случай негладких ограничений, в том числе диффференциальных. Здесь, напр., может быть использовано понятие субдифференциала выпуклой функции или его обобщения, когда выгуклость отсутствует (cM. [11], [12]).

Условия 1-го порядка, аналогичные принципу Понтрягина, известны для решений в классе обобщенных функций-мер (т. н. и н т е г р а ль н ы й п р и нц и п м а к с и му м а), для управляемых систем, описываемых дифференциальными уравнениями с отклоняющимся аргументом, дифференциальными уравнениями с частными производными, әволюционными уравнениями в банаховом пространстве, дифференциальными уравнениями на многообразиях, рекуррентными разностными уравнениями и т. д. (см. [1], [6], $[7],[13]-[16])$.

Из указанных необходимых условиї әкстремума в задаче оштимального управления вытекают известные необходимые условия 1-го порядка класспческого вариационного исчисления. В частности, в двухточечной краевой задаче для системы (2), (3), где $U$ - открытое множество, $J(x(\cdot), u(\cdot))$ - стандартный интегральный функционал, из принципа Понтрягина вытекает необходимое условие экстремума Вейерштрасса в классическом вариационном исчислении.

В теории оптимального управления развиваются методы получения необходимых условий высших порядков (особенно 2-го порядка) для неклассических вариационных задач (см. [19]). Интерес к условиям выспих порядков в значительной степени был связан с изучением вырожденных задач оптимального управления, приводящих к т. н. о с о б ы м у п р а в ле н ия м и не имеющих адекватных аналогов в классич. теории. Напр., в принципе Понтрягина фуункция $H(t, \psi$, $x, u)$ либо может приводить к целому семейству управлений, каждое из к-рых удовлетворяет принципу максимума, либо вообще не зависит от $u$ (тогда любое из допустимых значений $u$ удовлетворяет принципу Понтрягина). Әта ситуация оказалась весьма характерной для целого ряда прикладных задач управления в пространстве. В данном случае выделение оптимального управления уже требует перебора экстремалей 1-го порядка (т. н. э к с т р е м а л е ї̈ П о н т р я г и на) и применения к ним необходимых условий оптимальности 2-го или, в общем случае, более высокого порядка. Здесь разные по форме необходимые условия были получены путем использования специальных классов «неклассических" вариаций (напр., «связок» игольчатых вариаций и т. д.). Реализация особых управлений часто связана снова с использованием скользящих режимов (см. [17], [18]).

Теория достаточных условий оптимальности разработана в меньшей степени. Известны результаты, относящиеся к условиям локальной оптимальности и содержащие, в числе прочих требований, условия невырожденности системы в вариациях и ограничения на свойства гессиана правых частей, вычисленного вдоль исследуемой траектории для соответствующего обыкновенного дифференциального уравнения. Другая группа достаточных условий опирается на метод динамич. программирования и его связь с теорией принципа максимума (см. [8]). Имеются и формализмы, приводящие к достаточным условиям абсолютного минимума, основанные на идее распирения вариационных задач. Область их реальной применимости охватывает специальные классы задач с выпуклыми критериями и вырожденных задач оптимального управления (cM. [18]).

Полное решение задачи О. у. П. (необходимые и достаточные условия оптимальности) известно для линейных систем (1), когда рассматриваемые функционалы и ограничения на $u, x$ выпуклы (в ряде случаев здесь требуется выполнение нек-рых дополнительных условий). Привлечение идеї двойственности, используемых в выпуклом анализе, выявило особое әкстремальное свойство траекторий системы, описывающеӥ с о п р я ж е ны е п е р е ме н ны е принциіа максимума Понтрягина. Это позволило свести краевую задачу, возникающую при использовании необходимых условий общего вида, к решению более простой двойственной экстремальной задачи. В рамках названного подхода получила развитие теория линейіных систем с и м п у л ь с ны ы и у п р а в ле ни ям и, моделирующих объекты, подверженные мгновенным воздействиям (ударным, взрывным, импульсным), и формализуемых при помощи дифференциальных уравнений в обобщенных функциях соответствующих порядков сингулярности. Эффективное применение, особенно в теории игровых систем, нашел метод о б л а сте й дости и имости (см. [2], [3]).


\end{document}